\section{Introduction}\label{sec:intro}

A deployed model is often compared against a candidate while labels are scarce, delayed and
selectively collected. A common design labels only events on which the two models disagree, since
agreements carry no evidence about which is better, and promotes the candidate when an anytime-valid
certificate crosses a threshold \citep{karampatziakis2021offpolicy,choe2024comparing}. Promotion
changes the served model and carries a switching cost, so a false promotion is costly, and so is a
certificate that waits too long.

\paragraph{Which null.} The certificate is calibrated to a null, and published certificates use
different ones. Some control a conditional null: the conditional mean of the score increment given the
past is at most zero \citep{choe2024comparing}. Others assume an i.i.d.\ stream. A practitioner
usually cares about a marginal statement, namely that the challenger is not better on average. On
dependent streams the two differ. During an excursion the challenger \emph{is} conditionally ahead,
so a conditional-null certificate can promote on a stream whose marginal advantage is exactly zero
without violating its own guarantee. In a companion benchmark (P2 benchmark, confirmatory results;
\texttt{docs/p2\_benchmark/results/CONFIRMATORY\_RESULTS.md}) we measured this on exact real-data
marginal-null twins: certificate ports of the conditional-null family (``V1'') promoted at
$.098$ (99\% interval $[.091,.106]$), about $2\alpha$, and those promotions were conditionally
informed (mean remaining advantage $.0152$). % TODO: cite the benchmark paper once it has a reference.
A dependence-robust certificate (``V2'') removes this, but on synthetic dependent alternatives it
cost $+50.3$ per episode more than the cheapest V1 baseline in the same benchmark (confirmatory, 99\%
interval $[+47.0,+53.5]$, about $12\%$ of the cost $J$ defined in \cref{sec:setting}).

\paragraph{This paper.} We ask whether that price can be reduced without leaving a declared null
class. Our answer has a negative and a positive part. The negative part is that much of the price is
intrinsic: for a relaxed per-step class, the growth rate of any valid certificate is bounded by an
information rate that is several times smaller than the i.i.d.\ rate (\cref{prop:rate}), and
stationary-class methods may still beat it. The positive part is that a poorly designed robust
procedure pays more than this bound. Spending labels back-to-back wastes the budget, because after a
favourable outcome the envelope is still wide; spacing queries by the envelope is the direct repair.

\paragraph{Contributions.}
\begin{enumerate}
  \item \textbf{Envelope certificate.} A started mixture of betting factors tested against a
  per-step envelope $u_k$ on the conditional success probability under a declared dependence class,
  with \emph{predictable query spacing} and \emph{restart-safe conditioning}: the previous queried
  outcome is tracked over the whole stream, including outcomes a test restart discards
  (\cref{prop:validity}).
  \item \textbf{Price of robustness.} An information-rate characterisation (\cref{prop:rate}),
  stated for the \emph{relaxed} per-step class, with the caveat that procedures exploiting a fixed
  persistence on stationary data can beat it.
  \item \textbf{Two-stream design.} An envelope e-process on a flagged subsequence whose times are
  independent of the signs, combined with a universal-inference stream by a Bonferroni split
  (\cref{prop:twostream}).
  \item \textbf{Robustness value.} An anytime-valid dependence-robustness value, the largest
  persistence bound whose certificate has crossed (\cref{prop:rv}).
  \item \textbf{Payback certificate.} A promotion rule that bounds the probability of promoting when the
  forward net benefit does not cover the switching cost (\cref{prop:payback}).
  \item \textbf{Evaluation and benchmark correction.} A registered confirmation on the P2
  benchmark with exact marginal-null twins and 27 chronological replacement streams, with
  claimed-null validity and all-in cost reported together (\cref{sec:experiments}), including negative
  results. We also replace our earlier quad-endpoint MOPE port by an author-endpoint port whose
  interval width matches the authors' code (ratio $0.95$--$1.00$; the earlier port was $1.3$--$1.4\times$
  wider).
\end{enumerate}
\paragraph{What we find.} The registered confirmation gives a mixed result. Within the
dependence-robust class, NAD-2S is about $3\%$ cheaper than A5b and D5 on synthetic alternatives
(registered H1: $-11.31$, 99\% CI $[-13.79,-8.84]$; H3 on synthetic data: $-12.17$), keeps level
$\alpha$ in class ($.002$ promotion on both registered nulls) and passes the Tier B harm gate
(max 95\% UB $1.40$). On 27 weakly dependent real replacement streams it ties A5b and D5 ($+0.23$,
$+0.14$; registered H2 and H3 fail). Cheaper methods exist (author-endpoint MOPE and V1
front-loaders), but they promote on persistent marginal nulls ($.25$--$.32$ on the $\rho=.95$ null for
the author-endpoint MOPE rows; $.15$ pooled, 4 of 5 datasets above $\alpha$, on real twins in a
post-registration descriptive check). Those invalid promotions are economically harmless on the real
twins (mean decision harm $\le+0.8$). We therefore do \emph{not} claim state-of-the-art cost, and the
real-data economic motivation for robustness is weak; the characterisation of what robustness costs
is the main result.

We do not claim new probability theory. The results are elementary consequences of testing by
betting and Ville's inequality; the contribution is their combination in a budgeted, delayed-label
promotion rule and the measurement of what it costs.
